\chapter{Assignment 1: worked solutions and connections to the theory}

\begin{context}
After completing this assignment, you will be able to explain which continuity properties support
constructions by closure, decide in which algebra a spectrum is computed, and construct the commutative
subalgebra used to study a self-adjoint element. Each solution identifies the results it uses and how its
conclusion will be used later.
\end{context}

\section{Problem 1.1: continuity of multiplication and inversion}

Let $\A$ be a Banach algebra. In part (b), assume in addition that it has an identity.

\subsection*{1.1(a) Multiplication is jointly continuous}

\noindent\textbf{Statement.} If $x_n\to x$ and $y_n\to y$ in $\A$, prove that $x_ny_n\to xy$; equivalently, prove joint continuity of multiplication $\A\times\A\to\A$.

\begin{proof}[Solution]
Since $x_n\to x$, the sequence $(x_n)$ is bounded: there is $M$ with $\|x_n\|\le M$. Thus
\[
\|x_ny_n-xy\|
\le \|x_n\|\,\|y_n-y\|+\|x_n-x\|\,\|y\|
\le M\|y_n-y\|+\|y\|\|x_n-x\|\to0.
\]
Hence multiplication $\A\times\A\to\A$ is jointly continuous.
\end{proof}

\proofdiagram
{$x_n\to x$, $y_n\to y$}
{$\|ab\|\le\|a\|\|b\|$}
{Write $x_ny_n-xy=x_n(y_n-y)+(x_n-x)y$}
{$\|x_ny_n-xy\|\to0$}

\begin{cthejemplo}{(Product of convergent operator sequences)}
In $\mathcal B(H)$ we check Problem 1.1(a) with $T_n=T+n^{-1}I$ and $S_n=S+n^{-1}I$.
Then $T_nS_n-TS=n^{-1}(T+S)+n^{-2}I$, whence
$\|T_nS_n-TS\|\le n^{-1}(\|T\|+\|S\|)+n^{-2}\to0$.
This calculation shows that the product converges even when both factors vary simultaneously.
\end{cthejemplo}

\noindent\textbf{Dependencies.} The submultiplicative inequality in the definition of a Banach algebra and the boundedness of every convergent sequence. The estimate separates the variation of the two factors.

\noindent\textbf{How this problem is used later.} It allows passage to the closure of the algebras generated in Problem 1.3, passage to the limit in the normality relation in Problem 1.7, and control of polynomials in $x_n,x_n^*$ in Problem 2.8.

\subsection*{1.1(b) Inversion is continuous on $\GL(\A)$}

For this part, $\A$ is a unital Banach algebra and $\GL(\A)$ carries the norm inherited from
$\A$.

\noindent\textbf{Statement.} Prove that the map $\GL(\A)\to\GL(\A)$, $a\mapsto a^{-1}$, is continuous in the norm of $\A$.

\begin{proof}[Solution]
First, if $x\to1$ and $\|1-x\|<1$, the Neumann series gives
\[
 x^{-1}=\sum_{n=0}^\infty(1-x)^n,
 \qquad
 \|x^{-1}-1\|
 \le \sum_{n=1}^\infty\|1-x\|^n
 =\frac{\|1-x\|}{1-\|1-x\|}\to0.
\]
Therefore inversion is continuous at $1$. If $a_n\to a\in\GL(\A)$, then
$a^{-1}a_n\to1$ and, for all sufficiently large $n$, $a_n$ is invertible. Moreover,
\[
 a_n^{-1}=(a^{-1}a_n)^{-1}a^{-1}\longrightarrow a^{-1}.
\]
\end{proof}

\proofdiagram
{$a_n\to a\in\GL(\A)$}
{Continuity at $1$ by the Neumann series}
{$a_n=a(a^{-1}a_n)$ and $a^{-1}a_n\to1$}
{$a_n^{-1}=(a^{-1}a_n)^{-1}a^{-1}\to a^{-1}$}

\begin{cthejemplo}{(Inverse of a small perturbation)}
In $\mathcal B(H)$, we ask when a perturbation of the identity remains invertible.
If $E\in\mathcal B(H)$ satisfies $\|E\|<1$, the Neumann series gives
\[
 (I+E)^{-1}=\sum_{k=0}^{\infty}(-E)^k,
 \qquad \|(I+E)^{-1}-I\|\le\frac{\|E\|}{1-\|E\|}.
\]
The bound tends to zero with $\|E\|$: this is the concrete mechanism behind continuity in Problem 1.1(b).
\end{cthejemplo}

\noindent\textbf{Dependencies.} The Neumann series from Part 1 gives an inverse for $1-h$ if $\|h\|<1$; openness of $\GL(\A)$ allows the formula to be applied near any invertible element.

\noindent\textbf{How this problem is used later.} The proof of spectral permanence in Section 2 uses this continuity by going from $(b+i/n)^{-1}$ to $b^{-1}$ within a closed subalgebra.

\Needspace{14\baselineskip}
\section{Problem 1.2: commutation and inverses}

Let $\A$ be a unital complex algebra, let $a,b\in\A$, and assume that $a$ is invertible. In part (b),
assume in addition that $\A$ has an involution.

Recall that $x$ and $y$ \emph{$*$-commute} when
\[
 xy=yx,
 \qquad x^*y=yx^*.
\]

\subsection*{1.2(a) Commuting with $a$ is equivalent to commuting with $a^{-1}$}

\noindent\textbf{Statement.} Prove that $ab=ba$ if and only if $a^{-1}b=ba^{-1}$.

\begin{proof}[Solution]
If $ab=ba$, multiply first on the left and then on the right by $a^{-1}$:
\[
 b=a^{-1}ba \quad\Longrightarrow\quad ba^{-1}=a^{-1}b.
\]
Conversely, if $a^{-1}b=ba^{-1}$, multiply first on the left and then on the
right by $a$: from $b=aba^{-1}$ we obtain $ba=ab$. Thus $a$ commutes with $b$.
\end{proof}

\proofdiagram
{$ab=ba$, $a$ invertible}
{Multiply by $a^{-1}$ left and right}
{$a^{-1}b=ba^{-1}$}
{Equivalence, using $(a^{-1})^{-1}=a$}

\begin{cthejemplo}{(Commutation of diagonal matrices and their inverses)}
In $M_2(\mathbb C)$ we can see the equivalence of Problem 1.2(a) without resorting to limits.
Let's take $A=\operatorname{diag}(2,3)$ and $B=\operatorname{diag}(5,7)$.
As $A^{-1}=\operatorname{diag}(1/2,1/3)$,
\[
 AB=BA=\operatorname{diag}(10,21),\qquad
 A^{-1}B=BA^{-1}=\operatorname{diag}(5/2,7/3).
\]
Taking the inverse changes the entries of $A$ but preserves its commutation with $B$.
\end{cthejemplo}

\noindent\textbf{Dependencies.} Only associativity, the identity, and the existence of $a^{-1}$ are needed; no norm is involved.

\noindent\textbf{How this problem is used later.} In Problem 1.9 this identity is applied to $\lambda1-c$ to show that its inverse commutes with the chosen subalgebra.

\subsection*{1.2(b) $*$-commutation and inverses}

\noindent\textbf{Statement.} Prove that $a$ and $b$ $*$-commute if and only if $a^{-1}$ and $b$ $*$-commute.

\begin{proof}[Solution]
From $ab=ba$ and subpart (a), $a^{-1}b=ba^{-1}$. Additionally $a^*b=ba^*$; since $a^*$ is invertible,
$(a^*)^{-1}b=b(a^*)^{-1}$. Since $(a^*)^{-1}=(a^{-1})^*$, both conditions of
$*$-commutation hold. Conversely, if $a^{-1}$ and $b$ $*$-commute, apply the same argument
to the invertible element $a^{-1}$: its inverse $a$ and its adjoint $a^*$ commute with $b$.
\end{proof}

\Needspace{14\baselineskip}
\proofdiagram
{$a$ and $b$ $*$-commute}
{Subpart (a) applied to $a$ and $a^*$}
{$a^{-1}$ and $(a^{-1})^*=(a^*)^{-1}$ commute with $b$}
{$a^{-1}$ and $b$ $*$-commute}

\begin{cthejemplo}{(The inverse of a normal matrix)}
To check $*$-commutation in Problem 1.2(b), use the normal invertible matrix
$A=\operatorname{diag}(2,i)$ in $M_2(\mathbb C)$.
Then $A^{-1}=\operatorname{diag}(1/2,-i)$ and
\[
 A^{-1}(A^{-1})^*=(A^{-1})^*A^{-1}=\operatorname{diag}(1/4,1).
\]
Thus $A^{-1}$ remains normal: the key step is that its adjoint is the inverse of $A^*$.
\end{cthejemplo}

\noindent\textbf{Dependencies.} Problem 1.2(a) applies to both $a$ and $a^*$, along with $(a^{-1})^*=(a^*)^{-1}$. The definition of $*$-commutation sets the two relationships that need to be verified.

\noindent\textbf{How this problem is used later.} It controls inverses in the involutive version of Problem 1.9 and explains why the inverse of an invertible normal element is again normal.

\section{Problem 1.3: algebras generated by a set}

Let $\A$ be a Banach algebra and $S\subset\A$. We define
\[
\mathcal B(S)=\overline{\spn\{s_1\cdots s_k:k\ge1,\ s_j\in S\}}.
\]
If $\A$ has an involution, we put $S^*=\{s^*:s\in S\}$ and
$\mathcal B^*(S)=\mathcal B(S\cup S^*)$.

\subsection*{1.3(a) Minimality of $\mathcal B(S)$ and $\mathcal B^*(S)$}

\noindent\textbf{Statement.} Prove that $\mathcal B(S)$ is the smallest Banach subalgebra of $\A$ that contains $S$ and, when $\A$ is involutive, that $\mathcal B^*(S)$ is the minimal involutive Banach subalgebra containing $S$.

\begin{proof}[Solution]
The set of linear combinations of words in $S$ is a subalgebra; its closure is still a subalgebra by continuity of multiplication. Thus $\mathcal B(S)$ is a
Banach subalgebra containing $S$. If $C\subset\A$ is any Banach subalgebra with $S\subset C$,
then $C$ contains all the words $s_1\cdots s_k$, their linear combinations and, being
closed, its closure. Therefore $\mathcal B(S)\subset C$.

For $\mathcal B^*(S)$, apply the same argument to $S\cup S^*$. The generated algebra is stable under
$*$: the adjoint of a word is another word in $S\cup S^*$, and the involution is continuous. Hence it is the smallest involutive Banach subalgebra containing $S$.
\end{proof}

\pagebreak
\proofdiagram
{$S\subset\A$}
{Finite products, linear combinations and closure}
{Every closed subalgebra containing $S$ contains every word, linear combination, and limit}
{$\mathcal B(S)$ is the minimal Banach subalgebra containing $S$}

\begin{cthejemplo}{(Adding the unit changes the generated algebra)}
In $M_2(\mathbb C)$, the operator $T=\operatorname{diag}(0,2)$ distinguishes
$\mathcal B(\{T\})$ from the unital algebra generated by $T$.
As $T^k=\operatorname{diag}(0,2^k)$ for $k\ge1$,
\[
 \mathcal B(\{T\})=\{\operatorname{diag}(0,c):c\in\mathbb C\}.
\]
Adding $I$ gives all diagonal matrices. The distinction follows from the definition using words of positive length in Problem 1.3(a).
\end{cthejemplo}

\noindent\textbf{Dependencies.} The closure of linear combinations of words is complete by being closed in $\A$; Problem 1.1(a) allows products to pass to the closure. The continuity of $*$ gives the involutive case.

\noindent\textbf{How this problem is used later.} It constructs $C^*(a,1)$ concretely, the commutative subalgebra used in Section 2 for Gelfand--Naimark and functional calculus. The same generation operation appears in singular operator algebras.

\begin{investigacion}[title={Generation of a real operator algebra}]
To recognize Problem 1.3(a) in a research construction, consider the \emph{real} operator algebra in $\mathcal B(\mathbb R L^p(\Gamma,w))$ studied by Karlovich.
Its generators include $M_a$ with $a\in\SO^\diamond(\Gamma)$, $S_\Gamma$ and
real-linear conjugation $C$. For example, the word $M_aS_\Gamma C$ belongs to
the generated algebra, as do real linear combinations of words and their limits in
norm \cite[Eq.~(1.4), p.~3]{Karlovich2025}.
The generation operation is that of Problem 1.3(a) on real scalars. The presence of $C$
places this example outside the complex functional calculus of Section 2.
\end{investigacion}

\subsection*{1.3(b) commutation propagates to the generated algebra}

\noindent\textbf{Statement.} If $a\in\A$ commutes with each $s\in S$, prove commutation with every $x\in\mathcal B(S)$. In the involutive case, if $a$ $*$-commutes with each $s\in S$, prove the same conclusion for $\mathcal B^*(S)$ and for $*$-commutation.

\begin{proof}[Solution]
If $a$ commutes with each $s\in S$, then by induction it commutes with each word
$s_1\cdots s_k$. By linearity, it commutes with every linear combination of words. If
$x_n\to x$ and $ax_n=x_na$, continuity of multiplication gives $ax=x a$. Thus it commutes with all of
$\mathcal B(S)$.

In the involutive case, if $a$ $*$-commutes with each $s\in S$, then $a$ and $a^*$ commute with
$s$ and, taking adjoints, also with $s^*$. By the above argument, $a$ and $a^*$ commute with
all $\mathcal B^*(S)$, which is exactly to say that $a$ $*$-commutes with all its elements.
\end{proof}

\proofdiagram
{$as=sa$ for all $s\in S$}
{Commutation is preserved under products and sums}
{Pass to limits by continuity}
{$ax=xa$ for all $x\in\mathcal B(S)$}

\begin{cthejemplo}{(The commutant of diagonal matrices)}
In $M_2(\mathbb C)$ we want to follow the commutation from a generator to its algebra.
Let $D=\operatorname{diag}(1,2)$ and $A=\operatorname{diag}(3,4)$.
For any polynomial $p$,
\[
 Ap(D)=\operatorname{diag}(3p(1),4p(2))=p(D)A.
\]
The equality persists in polynomial limits in norm, as stated in Problem 1.3(b).
\end{cthejemplo}

\noindent\textbf{Dependencies.} The identity $as=sa$ is preserved when forming words and linear combinations; Problem 1.1(a) allows you to go to the limit. For $*$-commutation, the adjoint relations also apply.

\noindent\textbf{How this problem is used later.} Propagation of commutation justifies adding an inverse in Problem 1.9 without losing commutativity and helps analyze commutants of operator families.

\subsection*{1.3(c) If $S$ commutes pairwise, $\mathcal B(S)$ is commutative}

\noindent\textbf{Statement.} If $st=ts$ for all $s,t\in S$, prove that $\mathcal B(S)$ is commutative.

\begin{proof}[Solution]
Two words $s_1\cdots s_k$ and $t_1\cdots t_m$ can be rearranged factor by factor because all
generators commute. Therefore the words commute. By bilinearity the result extends to linear combinations of words, and by continuity to their closure.
\end{proof}

\proofdiagram
{The elements of $S$ commute pairwise}
{Reorder factors in two words}
{The words commute; then their linear combinations too}
{$\mathcal B(S)$ is commutative}

\begin{cthejemplo}{(Algebra of a nilpotent)}
The commutativity of an algebra generated by a single element does not require that it be normal.
In $M_2(\mathbb C)$, let's take $N=\begin{pmatrix}0&1\\0&0\end{pmatrix}$.
Since $N^2=0$, every word of positive length is $N$ or $0$, and
$\mathcal B(\{N\})=\mathbb C N$.
Two elements $cN,dN$ have product $0$ in either order. This illustrates Problem 1.3(c).
\end{cthejemplo}

\noindent\textbf{Dependencies.} The construction of Problem 1.3(a) is used: if the generators commute pairwise, any two words can be reordered. Problem 1.1(a) extends the identity to its limits.

\noindent\textbf{How this problem is used later.} It allows us to recognize the commutativity of $C^*(a,1)$ when $a$ is normal, a necessary step to represent that algebra as continuous functions in Section 2.

\subsection*{1.3(d) $*$-commutation pairwise}

\noindent\textbf{Statement.} In the involutive case, if each pair of elements of $S$ $*$-commutes, prove that $\mathcal B^*(S)$ is an involutive and commutative Banach subalgebra.

\begin{proof}[Solution]
Taking adjoints of $st=ts$ and $s^*t=ts^*$ also yields
$st^*=t^*s$ and $s^*t^*=t^*s^*$. Therefore all the elements of $S\cup S^*$ commute pairwise.
Applying part (c) to $S\cup S^*$ shows that $\mathcal B^*(S)$ is commutative; by construction,
it is an involutive Banach algebra.
\end{proof}

\proofdiagram
{Each pair $s,t\in S$ $*$-commutes}
{Take adjoints from relations}
{$S\cup S^*$ commutes pairwise}
{$\mathcal B^*(S)$ is commutative and involutive}

\begin{cthejemplo}{(A normal element generates a commutative involutive algebra)}
In $M_2(\mathbb C)$, the matrix $A=\operatorname{diag}(1,i)$ shows why
normality allows using functional calculus.
Its adjoint is $A^*=\operatorname{diag}(1,-i)$, so $AA^*=A^*A=I$.
Thus words in $A$, $A^*$, and $I$ are diagonal and commute, as does their closure.
By Problem 1.3(d), $C^*(A,I)$ is commutative, as required to apply Section 2.
\end{cthejemplo}

\noindent\textbf{Dependencies.} The definition of $*$-commutation applied to each pair of generators gives commutation in $S\cup S^*$; then Problem 1.3(c) is applied.

\noindent\textbf{How this problem is used later.} If $a$ is normal, applying this part to $S=\{a,1\}$ gives the commutative $C^*$-subalgebra $C^*(a,1)$, starting point of continuous functional calculus.

\section{Problem 1.4: the quotient by an ideal is a Banach algebra}

Let $\A$ be a Banach algebra and $J\subset\A$ be a closed two-sided ideal.

\noindent\textbf{Statement.} Prove that $(x+J)(y+J)=xy+J$ is well defined and that $\A/J$, with $\|x+J\|=\inf_{j\in J}\|x-j\|$, is a Banach algebra.

\begin{proof}[Solution]
If $x'=x+j_1$ and $y'=y+j_2$ with $j_1,j_2\in J$, then
\[
 x'y'-xy=xj_2+j_1y+j_1j_2\in J,
\]
so $(x+J)(y+J)=xy+J$ is well defined.

For the norm, take $j_1,j_2\in J$: 
\[
(x-j_1)(y-j_2)=xy-\underbrace{(xj_2+j_1y-j_1j_2)}_{\in J},
\]
so
\[
\|xy+J\|\le\|x-j_1\|\,\|y-j_2\|.
\]
Taking the infimum over $j_1$ and $j_2$ gives
$\|(x+J)(y+J)\|\le\|x+J\|\|y+J\|$. Finally, since $J$ is closed, the quotient space
$\A/J$ is complete. This is a Banach algebra.
\end{proof}

\proofdiagrambranch
{$J$ closed two-sided ideal}
{$\|x+J\|=\inf_{b\in J}\|x-b\|$}
{Check independence of representatives and submultiplicativity}
{Use that the quotient by a closed subspace is Banach}
{$\A/J$ is a Banach algebra}

\begin{cthejemplo}{(A class in the Calkin algebra)}
To see what the quotient in Problem 1.4 identifies, work in $X=\ell^2(\mathbb N)$.
Let $P$ be the rank one projection onto $\mathbb C e_1$; then $P\in\K(X)$ and
\[
 (I+P)+\K(X)=I+\K(X).
\]
Two different operators represent the same class because their difference is compact.
This is the basic equality that allows us to work in the Calkin algebra.
\end{cthejemplo}

\noindent\textbf{Dependencies.} A two-sided ideal ensures that the product of classes does not depend on representatives; closedness of $J$ ensures completeness of the quotient. The quotient norm and submultiplicativity in $\A$ are used.

\noindent\textbf{How this problem is used later.} By taking $\A=\mathcal B(X)$ and $J=\K(X)$ we obtain the Calkin algebra. The equivalence between Fredholm and invertibility in that quotient is Atkinson's theorem, an additional result that this exercise allows us to formulate but does not prove.

\begin{investigacion}[title={Fredholm operators in the Calkin algebra}]
To place Problem 1.4 in an application, let $X$ be an infinite-dimensional Banach space and
$\pi:\mathcal B(X)\to\mathcal B(X)/\K(X)$ the quotient map.
The class $\pi(A)$ is invertible exactly when there exists $B\in\mathcal B(X)$ with
$AB-I,BA-I\in\K(X)$; both equalities follow by multiplying classes in the quotient.
Atkinson's theorem adds that this condition is equivalent to $A$ being Fredholm: range
closed, finite-dimensional kernel, and finite-dimensional quotient $X/\operatorname{Ran}(A)$
\cite[p.~2]{Karlovich2025}. Problem 1.4 provides the algebra in which this equivalence is formulated; Atkinson's theorem is an additional result.
\end{investigacion}

\section{Problem 1.5: \texorpdfstring{$C_0(X)$}{C0(X)} is a Banach subalgebra of \texorpdfstring{$C_b(X)$}{Cb(X)}}

Let $X$ be a locally compact and Hausdorff space.

\noindent\textbf{Statement.} Define
\begin{multline*}
C_0(X)=\{f\in C_b(X):\text{ for every }\varepsilon>0\text{ there is a compact }K\subset X\\
\text{ such that }|f(x)|<\varepsilon\text{ outside }K\}.
\end{multline*}
Prove that it is a Banach subalgebra of $C_b(X)$.

\begin{proof}[Solution]
Any continuous function that vanishes at infinity is bounded: for $\varepsilon=1$ it is less than $1$ outside a compact set, and over that compact it reaches a finite maximum.
Also, if $f,h\in C_0(X)$, we choose compact $K_f,K_h$ outside which
$|f|,|h|<\varepsilon/2$; outside $K_f\cup K_h$ we have $|f+h|<\varepsilon$.
Multiplication by scalars is verified by choosing the threshold
$\varepsilon/(1+|c|)$ for $cf$. Thus $C_0(X)$ is a linear subspace of $C_b(X)$.

If $f\in C_0(X)$ and $g\in C_b(X)$, given $\varepsilon>0$ we choose a compact $K$ such that
$|f(x)|<\varepsilon/(1+\|g\|_\infty)$ outside $K$. So
$|f(x)g(x)|<\varepsilon$ outside $K$, so $fg\in C_0(X)$. This even shows that $C_0(X)$
is an ideal of $C_b(X)$.

To show closedness, let $f_n\in C_0(X)$ and $f_n\to f$ uniformly. Given $\varepsilon>0$,
we choose $n$ with $\|f-f_n\|_\infty<\varepsilon/2$ and a compact $K$ with
$|f_n(x)|<\varepsilon/2$ outside $K$. Then $|f(x)|<\varepsilon$ outside $K$.
Therefore $f\in C_0(X)$. A closed subspace of a Banach space is itself a Banach space.

Pointwise conjugation preserves $C_0(X)$ and
$\|f^*f\|_\infty=\sup_x|f(x)|^2=\|f\|_\infty^2$; therefore $C_0(X)$ is also a
$C^*$-algebra, although it has no unit when $X$ is not compact.
\end{proof}

\proofdiagram
{$f\in C_0(X)$, $g\in C_b(X)$}
{Definition of vanishing at infinity}
{$fg$ goes small out of a compact; uniform limits preserve the property}
{$C_0(X)$ is a closed ideal of $C_b(X)$, then Banach}

\begin{cthejemplo}{(A function that vanishes at infinity on the real line)}
In $X=\mathbb R$ we want to compare a vanishing function with the unit of $C_b(X)$.
For $f(x)=1/(1+x^2)$ and $\varepsilon>0$, if $|x|>\varepsilon^{-1/2}$ then
$|f(x)|<\varepsilon$; therefore $f\in C_0(\mathbb R)$.
The constant function $1$ does not satisfy that condition for $\varepsilon=1/2$ outside of any
compact. Problem 1.5 places $f$ in the closed ideal $C_0(\mathbb R)$ and leaves the unit outside of it.
\end{cthejemplo}

\noindent\textbf{Dependencies.} The definition of vanishing at infinity allows a compact to be chosen for each tolerance. The $C_b(X)$ uniform norm controls both products and limits.

\noindent\textbf{How this problem is used later.} If $X$ is not compact, $C_0(X)$ is the non-unital commutative model; for compact $X$, it coincides with $C(X)$. The general classification of $C_0(X)$ requires the non-unital version of Gelfand--Naimark, while Section 2 demonstrates the unital version.

\section{Problem 1.6: spectrum and involution}

Let $\A$ be a unital involutive Banach algebra and $x\in\A$.

\noindent\textbf{Statement.} For each $x\in\A$, prove $\spec(x^*)=\overline{\spec(x)}$.

\begin{proof}[Solution]
An element $a$ is invertible if and only if $a^*$ is, since from $ab=ba=1$ we obtain
$b^*a^*=a^*b^*=1$. Now
\[
(\lambda1-x^*)^*=\overline\lambda1-x.
\]
Therefore
\[
\lambda1-x^*\notin\GL(\A)
\Longleftrightarrow
\overline\lambda1-x\notin\GL(\A),
\]
which is equivalent to
\[
 \spec(x^*)=\overline{\spec(x)}.
\]
\end{proof}

\proofdiagram
{$\lambda1-x^*$}
{involution preserves invertibility: $a^{-1}\mapsto(a^*)^{-1}$}
{$(\lambda1-x^*)^*=\overline\lambda1-x$}
{$\lambda\in\spec(x^*)\Longleftrightarrow\overline\lambda\in\spec(x)$}

\begin{cthejemplo}{(Spectrum of a non-normal adjoint)}
The symmetry of Problem 1.6 does not require normality. Let's consider
$A=\begin{pmatrix}1&1\\0&i\end{pmatrix}$ in $M_2(\mathbb C)$.
Since $A$ is triangular, $\spec(A)=\{1,i\}$; its adjoint is lower triangular with diagonal
$(1,-i)$, so $\spec(A^*)=\{1,-i\}=\overline{\spec(A)}$.
Conjugation of the spectrum follows from the involution, even though $A$ does not commute with $A^*$.
\end{cthejemplo}

\noindent\textbf{Dependencies.} The spectrum definition of Part 1 and the identity $(a^*)^{-1}=(a^{-1})^*$ allow invertibility to be carried through the involution.

\noindent\textbf{How this problem is used later.} It shows a spectral symmetry that holds before imposing the identity $C^*$. To conclude that the entire spectrum of a self-adjoint is real, Section 2 adds arguments proper to $C^*$-algebras.

\section{Problem 1.7: normal elements form a closed set}

Let $\A$ be a $C^*$-algebra and $x_n\to x$ in norm, with each $x_n$ being normal.

\noindent\textbf{Statement.} Prove that the set $\{x\in\A:x^*x=xx^*\}$ of normal elements is closed in norm.

\begin{proof}[Solution]
As $\|x_n^*-x^*\|=\|x_n-x\|\to0$ and multiplication is jointly continuous,
\[
 x_n^*x_n\to x^*x,
 \qquad
 x_nx_n^*\to xx^*.
\]
But both terms are the same for each $n$. By uniqueness of limits, $x^*x=xx^*$.
\end{proof}

\proofdiagram
{$x_n\to x$, $x_n^*x_n=x_nx_n^*$}
{The involution is isometric; jointly continuous product}
{$x_n^*\to x^*$ and both products converge}
{$x^*x=xx^*$}

\begin{cthejemplo}{(Limit of normal matrices)}
In $M_2(\mathbb C)$ we can observe the closure of the set of normals without changing dimension.
Let $A_n=\operatorname{diag}(1+1/n,i-1/n)$; each $A_n$ is normal and
$\|A_n-\operatorname{diag}(1,i)\|=1/n$.
The limit $A=\operatorname{diag}(1,i)$ also commutes with its adjoint.
This is a concrete instance of the passage to the limit proved in Problem 1.7.
\end{cthejemplo}

\noindent\textbf{Dependencies.} The isometric involution gives $x_n^*\to x^*$ and Problem 1.1(a) allows going to the limit in the two products of the normality equation.

\noindent\textbf{How this problem is used later.} It is the first step of Problem 2.8: the limit $x$ is still normal, so $f(x)$ is defined by continuous functional calculus.

\section{Problem 1.8: spectral inclusion under homomorphisms}

Let $\A,\B$ be unital Banach algebras and $\psi:\A\to\B$ be a unital complex homomorphism.

\noindent\textbf{Statement.} For each $x\in\A$, prove $\spec_\B(\psi(x))\subset\spec_\A(x)$ and $r_\B(\psi(x))\le r_\A(x)$.

\begin{proof}[Solution]
If $\lambda\notin\spec_\A(x)$, then $(\lambda1-x)^{-1}$ exists. Applying $\psi$,
\[
 \psi(\lambda1-x)\,\psi((\lambda1-x)^{-1})
 =\psi((\lambda1-x)^{-1})\,\psi(\lambda1-x)=1_\B,
\]
so $\lambda1-\psi(x)$ is invertible. This proves
$\spec_\B(\psi(x))\subset\spec_\A(x)$. Taking maxima of $|\lambda|$ we obtain
$r(\psi(x))\le r(x)$.
\end{proof}

\proofdiagram
{$\psi:\A\to\B$ unital homomorphism}
{Unital homomorphisms send inverses to inverses}
{$\lambda1-\psi(x)=\psi(\lambda1-x)$}
{$\spec(\psi(x))\subset\spec(x)$ and $r(\psi(x))\le r(x)$}

\begin{cthejemplo}{(Strict spectral inclusion in a quotient)}
To show that the inclusion of Problem 1.8 can be strict, let's take
$\A=\mathbb C\oplus\mathbb C$ with supremum norm and $J=\{0\}\oplus\mathbb C$.
The class of $a=(1,2)$ in $\A/J\cong\mathbb C$ is $1$.
As $\lambda1-a=(\lambda-1,\lambda-2)$ is invertible exactly when
$\lambda\notin\{1,2\}$,
\[
 \spec_{\A/J}(a+J)=\{1\}\subsetneq\{1,2\}=\spec_\A(a).
\]
The quotient may lose spectral points, according to Problem 1.8.
\end{cthejemplo}

\noindent\textbf{Dependencies.} A unital homomorphism preserves the inverse of $\lambda1-x$ when it exists; the definition of spectrum converts this algebraic observation into an inclusion.

\noindent\textbf{How this problem is used later.} In Section 2 we apply inclusion to $a^*a$ to prove that unital $*$-homomorphisms are contractive.

\section{Problem 1.9: a spectrum-preserving commutative subalgebra}

Let $\A$ be a unital Banach algebra and $S\subset\A$ be a family that commutes pairwise. For the
involutive version, we further assume that $\A$ is involutive and the elements of $S$
$*$-commute pairwise.

\noindent\textbf{Statement.} Prove that there exists a unital and commutative Banach subalgebra $C\subset\A$ that contains $S$ and satisfies $\spec_C(s)=\spec_\A(s)$ for every $s\in S$. In the involutive case with $*$-commutation pairwise, require in addition that $C$ be involutive.

\begin{proof}[Solution]
Let $\mathscr F$ be the family of unital commutative Banach subalgebras $C\subset\A$ containing $S$, ordered by inclusion. It is non-empty: it contains the closure of the unital algebra generated by
$S$. If $\{C_i\}$ is a chain, the closure of $\bigcup_i C_i$ is again a unital commutative Banach subalgebra containing $S$. By Zorn's lemma there exists a maximal element $C$.

We assert that $C$ is inverse-closed in $\A$. If $c\in C$ is invertible in $\A$, then
from $cd=dc$ for $d\in C$ it follows that $c^{-1}d=dc^{-1}$. Therefore the closed unital subalgebra
generated by $C\cup\{c^{-1}\}$ is still commutative. Maximality forces
$c^{-1}\in C$. Consequently, for $s\in S$,
\[
\lambda1-s\in\GL(\A)\Longleftrightarrow\lambda1-s\in\GL(C),
\]
and the spectra match.

In the involutive case, we start from the family of involutive, commutative, unital Banach subalgebras containing $S$. The assumption of $*$-commutation guarantees that the family is non-empty by
Problem 1.3(d), applied to $S\cup\{1\}$. If $c\in C$ is invertible in $\A$, both $c^{-1}$ and $(c^{-1})^*$ commute with
$C$ and each other: $c,c^*\in C$ commute and their inverses as well. They can be added while keeping
commutativity and involution. Maximality again forces $c^{-1}\in C$.
\end{proof}

\proofdiagrambranch
{$S$ commutes pairwise}
{Family of closed, commutative, unital subalgebras containing $S$}
{Zorn produces a maximal; the inverses of its elements also commute with it}
{Maximality forces the inverses to belong to the subalgebra}
{$\spec_\A(s)=\spec_C(s)$ for $s\in S$}

\begin{cthejemplo}{(A spectrally faithful subalgebra of matrices)}
In $M_2(\mathbb C)$, the self-adjoint $A=\operatorname{diag}(1,2)$ allows building
explicitly the subalgebra of Problem 1.9.
Let's take $C$ as all diagonal matrices. The inverse of
$\lambda I-A=\operatorname{diag}(\lambda-1,\lambda-2)$, when it exists, still belongs to $C$:
\[
 (\lambda I-A)^{-1}=\operatorname{diag}((\lambda-1)^{-1},(\lambda-2)^{-1}).
\]
Therefore $\spec_C(A)=\{1,2\}=\spec_{M_2(\mathbb C)}(A)$; the spectrum is conserved because
the necessary inverses belong to $C$.
\end{cthejemplo}

\noindent\textbf{Dependencies.} Problem 1.3 is used to generate a commutative subalgebra, Problem 1.2 to commute with inverses and Zorn's lemma to choose a maximal subalgebra. The involutive version requires $*$-commutation.

\noindent\textbf{How this problem is used later.} Section 2 uses this reduction to study a self-adjoint using characters without changing its spectrum. In the $C^*$ environment, spectral permanence then offers a more direct path to $C^*(a,1)$.

\pagebreak
\section{Problem 1.10: spectrum of a continuous function}

Let $X$ be compact Hausdorff and $f\in C(X)$. Then we will compare with $C_b(X)$ for a locally compact non-compact space.

\noindent\textbf{Statement.} Prove that $\spec_{C(X)}(f)=f(X)$ for each $f\in C(X)$ and decide, using an example, whether the equality still holds in $C_b(X)$ for $X$ only locally compact.

\begin{proof}[Solution]
If $\lambda\notin f(X)$, the function $\lambda-f$ has no zeros. Since $X$ is compact,
$|\lambda-f|$ has positive minimum; therefore
$(\lambda-f)^{-1}\in C(X)$ and $\lambda\notin\spec(f)$. If $\lambda=f(x_0)$, then
$\lambda-f$ vanishes at $x_0$ and cannot be invertible. Therefore $\spec(f)=f(X)$.

If $X$ is only locally compact and we use $C_b(X)$, the statement may fail because a nonvanishing function can approach zero without reaching zero. For example, in $X=(0,1)$, $f(x)=x$ has
$f(X)=(0,1)$, but
\[
 \spec_{C_b((0,1))}(f)=[0,1].
\]
Indeed, for $\lambda\in(0,1)$ the function $\lambda-f$ vanishes; for $\lambda=0,1$,
their reciprocals are unbounded. If $\lambda\notin[0,1]$, however,
$|\lambda-f(x)|\ge\operatorname{dist}(\lambda,[0,1])>0$ for all $x$, so its reciprocal belongs to $C_b((0,1))$. This proves the announced equality.
\end{proof}

\proofdiagram
{$f\in C(X)$, $X$ compact}
{Pointwise invertibility in an algebra of functions}
{$\lambda1-f$ invertible $\Longleftrightarrow$ $\lambda-f$ does not vanish}
{$\spec(f)=f(X)$}

\begin{cthejemplo}{(A spectrum shaped like a circle)}
To read Problem 1.10 geometrically, consider $f(z)=2+z$ in $C(\mathbb T)$.
Function $\lambda-f$ vanishes on $\mathbb T$ exactly when
$\lambda=2+z$ for some $|z|=1$; outside those values it has a continuous reciprocal.
Therefore
\[
 \spec_{C(\mathbb T)}(f)=\{\lambda\in\mathbb C:|\lambda-2|=1\}.
\]
The spectrum is the image of $f$, a circle with center $2$ and radius $1$.
\end{cthejemplo}

\noindent\textbf{Dependencies.} In $C(X)$, $\lambda1-f$ is invertible exactly when $\lambda-f$ does not vanish; compactness of $X$ ensures that its reciprocal is continuous and bounded.

\noindent\textbf{How this problem is used later.} Identifies the spectrum with a range of values in the Gelfand--Naimark representation and is the specific step that Problem 2.4 uses to demonstrate spectral mapping.

\pagebreak
\section{Problem 1.11: \texorpdfstring{$\mathcal B(H)$}{B(H)} and the unilateral shift}

Let $H$ be a Hilbert space and $T\in\mathcal B(H)$. Subsections (b)--(d) change the assumptions
on $H$ as stated in each part.

\subsection*{1.11(a) $\mathcal B(H)$ is a $C^*$-algebra}

\noindent\textbf{Statement.} Prove that $\mathcal B(H)$, with the operator norm and the involution $T\mapsto T^*$, is a $C^*$-algebra.

\begin{proof}[Solution]
Part 1 already established that bounded operators on a Banach space form a
Banach algebra; therefore $\mathcal B(H)$ is complete. The adjoint rules give an involution.
The norm of the adjoint coincides with that of $T$: by the characterization of the norm of a vector of
a Hilbert space using inner products,
\[
 \|T^*\|=\sup_{\|\eta\|=\|\xi\|=1}|\langle T^*\eta,\xi\rangle|
 =\sup_{\|\eta\|=\|\xi\|=1}|\langle\eta,T\xi\rangle|=\|T\|.
\]
Submultiplicativity then gives $\|T^*T\|\le\|T^*\|\|T\|=\|T\|^2$.
For the reverse inequality, if $\|\xi\|=1$,
\[
\|T\xi\|^2=\langle T^*T\xi,\xi\rangle\le\|T^*T\|.
\]
Taking the supremum over $\xi$ gives $\|T\|^2\le\|T^*T\|$.
\end{proof}

\proofdiagram
{$T\in\mathcal B(H)$}
{adjoint properties}
{$\|T^*T\|\le\|T\|^2$ and $\|T\xi\|^2=\langle T^*T\xi,\xi\rangle$}
{$\|T^*T\|=\|T\|^2$}

\begin{cthejemplo}{(The identity $C^*$ in a diagonal matrix)}
In $H=\mathbb C^2$ we want to check the equality of norms that distinguishes
$\mathcal B(H)$ as $C^*$-algebra.
For $T=\operatorname{diag}(2,i)$, there is
$T^*T=\operatorname{diag}(4,1)$, then $\|T^*T\|=4$ and $\|T\|=2$.
Thus $\|T^*T\|=\|T\|^2$, the identity verified generally in Problem 1.11(a).
\end{cthejemplo}

\noindent\textbf{Dependencies.} Part 1 presents $\mathcal B(H)$ as a Banach algebra; the properties of the adjoint and the operator norm allow the $C^*$-identity to be verified.

\noindent\textbf{How this problem is used later.} Provides the concrete model from Section 2 and the operator space in which Problem 2.1 isometrically represents $C_b(X)$.

\subsection*{1.11(b) In finite dimension, spectrum = eigenvalues}

\noindent\textbf{Statement.} If $\dim H<\infty$, prove that $\spec(T)$ matches the set of eigenvalues of $T$.

\begin{proof}[Solution]
In finite dimension a linear operator is invertible if and only if its kernel is zero. So
$\lambda I-T$ is noninvertible exactly when $\xi\neq0$ exists with
$(\lambda I-T)\xi=0$, that is, $T\xi=\lambda\xi$.
\end{proof}

\proofdiagram
{$\dim H<\infty$}
{In finite dimension: injective $\Longleftrightarrow$ surjective}
{$\lambda I-T$ non-invertible $\Longleftrightarrow\Ker(\lambda I-T)\neq0$}
{$\lambda\in\spec(T)\Longleftrightarrow\lambda$ eigenvalue}

\begin{cthejemplo}{(Eigenvalues of a triangular matrix)}
In dimension two, the determinant allows the spectrum and
eigenvectors of Problem 1.11(b).
For $T=\begin{pmatrix}1&1\\0&2\end{pmatrix}$,
$\det(\lambda I-T)=(\lambda-1)(\lambda-2)$.
Also $Te_1=e_1$ and $T(1,1)=(2,2)$, so $1$ and $2$ are eigenvalues.
There are no other spectral points: in finite dimension the singularity of $\lambda I-T$
matches the existence of an eigenvector.
\end{cthejemplo}

\noindent\textbf{Dependencies.} In finite dimension, injectivity and surjectivity of $\lambda I-T$ are equivalent by rank and nullity. The definition of spectrum from Part 1 applies.

\noindent\textbf{How this problem is used later.} Allows you to interpret the matrix examples in Section 2 as spectrum calculations using eigenvalues; Problem 1.11(c)--(d) show what changes in infinite dimension.

\subsection*{1.11(c) Approximate eigenvectors produce spectral points}

\noindent\textbf{Statement.} If $\lambda\in\mathbb C$ and there exist unit vectors $\xi_n\in H$ such that $\|(\lambda I-T)\xi_n\|\to0$, prove that $\lambda\in\spec(T)$.

\begin{proof}[Solution]
If $\lambda\notin\spec(T)$, then $R=(\lambda I-T)^{-1}$ is bounded. Therefore
\[
1=\|\xi_n\|=\|R(\lambda I-T)\xi_n\|
\le\|R\|\,\|(\lambda I-T)\xi_n\|\to0,
\]
contradiction. Then $\lambda\in\spec(T)$.
\end{proof}

\proofdiagram
{$\|\xi_n\|=1$, $\|(\lambda I-T)\xi_n\|\to0$}
{Assume $\lambda I-T$ invertible}
{$\xi_n=(\lambda I-T)^{-1}(\lambda I-T)\xi_n$}
{$1\le\|(\lambda I-T)^{-1}\|\|(\lambda I-T)\xi_n\|\to0$, contradiction}

\begin{cthejemplo}{(Approximate spectrum without eigenvector)}
In $\ell^2(\mathbb N)$ we want to see that an approximate spectral point can lack
eigenvector. Let's define $T(\xi_n)=(\xi_n/n)$ and take $\lambda=0$.
Each $e_n$ is a unit vector and $\|Te_n\|=1/n\to0$, but $T\xi=0$ implies
$\xi_n/n=0$ for all $n$ i.e. $\xi=0$.
By Problem 1.11(c), $0\in\spec(T)$; the calculation distinguishes approximate spectrum from point spectrum.
\end{cthejemplo}

\noindent\textbf{Dependencies.} If $\lambda I-T$ is invertible, its inverse is bounded. The norm estimate of the composition contradicts the existence of almost null unit vectors.

\noindent\textbf{How this problem is used later.} Defines the approximate point spectrum and prepares the comparison with the unilateral shift of Problem 1.11(d), whose spectral zero is not an approximate eigenvalue.

\subsection*{1.11(d) Unilateral shift shows the difference}

\noindent\textbf{Statement.} Let $S(x_1,x_2,\ldots)=(0,x_1,x_2,\ldots)$ be
$\ell^2(\mathbb N)$. Prove that $0\in\spec(S)$ and that there are no unit vectors
$\xi_n$ with $\|S\xi_n\|\to0$.

\begin{proof}[Solution]
$S$ is not surjective: no vector with non-zero first coordinate is in its range. So $S$ is not invertible and $0\in\spec(S)$. But
\[
 \|S\xi\|^2=0+\sum_{n\ge1}|\xi_n|^2=\|\xi\|^2.
\]
Therefore, for every unit vector $\xi$, $\|S\xi\|=1$; a sequence of unit vectors cannot exist with
$\|S\xi_n\|\to0$.
\end{proof}

\proofdiagram
{$S(x_1,x_2,\dots)=(0,x_1,x_2,\dots)$}
{$S$ is isometry but is not surjective}
{$0\in\spec(S)$; however $\|S\xi\|=\|\xi\|$}
{No unit vectors with $\|S\xi_n\|\to0$}

\begin{cthejemplo}{(An isometry that is not surjective)}
To distinguish conservation of norm and invertibility in Problem 1.11(d), let's take the first
$\ell^2(\mathbb N)$ base vector.
The shift satisfies $Se_1=e_2$ and $\|Se_1\|=\|e_1\|=1$, but no vector can be mapped to $e_1$, because the first coordinate of $S\xi$ is always zero.
Isometry is not surjective; that is why $0$ belongs to its spectrum without having unit vectors that are nearly annihilated.
\end{cthejemplo}

\noindent\textbf{Dependencies.} The unilateral shift is an isometry, but its range omits vectors with a non-zero first coordinate. The definition of spectrum detects this lack of surjectivity.

\noindent\textbf{How this problem is used later.} Problem 2.10 will determine the entire spectral disk using eigenvectors of adjoint. Isometry without surjectivity explains why $0$ can be spectral without being an approximate eigenvalue.

\begin{investigacion}[title={Two shifts in different spaces}]
To recognize the scope of Problem 1.11(d), let's compare the shift $S$ of $\ell^2(\mathbb N)$ with
$V_\alpha$ in Haseman's formulation. The first one satisfies $Se_1=e_2$; the second
composes boundary functions, $(V_\alpha f)(t)=f(\alpha(t))$
\cite[pp.~463--464, ecs.~(1.1)--(1.2)]{Karlovich2017}.
Both move data, but act in different spaces and geometries. Problem 1.11(d)
determines $0\in\spec(S)$; it does not determine by analogy any point of $\spec(V_\alpha)$.
\end{investigacion}

\section{Problem 1.12: the disk algebra and failure of spectral permanence in general}

Let $\mathbb D=\{z:|z|<1\}$, $\mathbb T=\{z:|z|=1\}$ and
\[
 A(\mathbb D)=\{f\in C(\overline{\mathbb D}): f\text{ is holomorphic on }\mathbb D\}.
\]

\subsection*{1.12(a) $A(\mathbb D)$ is a Banach algebra}

\noindent\textbf{Statement.} Prove that $A(\mathbb D)$ is a unital Banach algebra with the norm of the supremum in $\overline{\mathbb D}$.

\begin{proof}[Solution]
Let $(f_n)$ be Cauchy in the supremum norm. It converges uniformly on $\overline{\mathbb D}$ to a
continuous function $f$. In each $K\subset\mathbb D$ compact, the convergence is uniform; by
Weierstrass' theorem, $f$ is holomorphic in $\mathbb D$. So $f\in A(\mathbb D)$ and the space is
complete. The sums and products of holomorphic functions are holomorphic, and the constant function
$1$ belongs to $A(\mathbb D)$; operations and norm are inherited from
$C(\overline{\mathbb D})$.
\end{proof}

\proofdiagram
{Cauchy sequence in supremum norm}
{$C(\overline{\mathbb D})$ is Banach; Weierstrass theorem for holomorphic limits}
{The uniform limit is continuous and the limit in interior compacts is holomorphic}
{$A(\mathbb D)$ is closed, unital and Banach}

\begin{cthejemplo}{(Polynomial approximation in disk algebra)}
In $A(\mathbb D)$ we are interested in approximating holomorphic functions without leaving the uniform norm.
For $f\in A(\mathbb D)$, the dilations $f_r(z)=f(rz)$ converge uniformly to $f$
in $\overline{\mathbb D}$ when $r\uparrow1$, by uniform continuity of $f$.
Each $f_r$ is holomorphic on a disk of radius $1/r>1$, so its Taylor polynomials
converge uniformly on $\overline{\mathbb D}$.
This explains why the polynomials form a dense subalgebra of the Banach algebra in Problem 1.12(a).
\end{cthejemplo}

\noindent\textbf{Dependencies.} The completeness of $C(\overline{\mathbb D})$ and the Weierstrass theorem on local uniform limits of holomorphic functions guarantee that $A(\mathbb D)$ is closed.

\noindent\textbf{How this problem is used later.} It makes it possible to treat the boundary restriction as an inclusion of Banach algebras in Problem 1.12(b) and to compare spectra in Problem 1.12(c).

\subsection*{1.12(b) The boundary restriction is isometric}

\noindent\textbf{Statement.} Prove that the restriction $f\mapsto f|_{\mathbb T}$ is a unital isometric homomorphism from $A(\mathbb D)$ into $C(\mathbb T)$.

\begin{proof}[Solution]
By the maximum modulus principle,
\[
 \|f\|_{C(\overline{\mathbb D})}=\max_{|z|=1}|f(z)|=\|f|_{\mathbb T}\|_{C(\mathbb T)}.
\]
Thus the restriction is isometric and injective. For $f,g\in A(\mathbb D)$ and $z\in\mathbb T$,
\[
 (f+g)|_{\mathbb T}(z)=f(z)+g(z),\qquad
 (fg)|_{\mathbb T}(z)=f(z)g(z),\qquad
 1|_{\mathbb T}(z)=1.
\]
Therefore it is a unital homomorphism of Banach algebras.
\end{proof}

\proofdiagram
{$f\in A(\mathbb D)$}
{Maximum modulus principle for holomorphic functions}
{$\max_{\overline{\mathbb D}}|f|=\max_{\mathbb T}|f|$}
{$f\mapsto f|_{\mathbb T}$ is an isometric embedding}

\begin{cthejemplo}{(The disk norm is attained on the boundary)}
In $A(\mathbb D)$ we check Problem 1.12(b) with $f(z)=1+z^2$.
For $|z|\le1$, $|f(z)|\le2$, and equality is achieved in $z=\pm1\in\mathbb T$.
So $\|f\|_{C(\overline{\mathbb D})}=\|f|_{\mathbb T}\|_{C(\mathbb T)}=2$:
The boundary already determines the norm of this function.
\end{cthejemplo}

\noindent\textbf{Dependencies.} The maximum modulus principle identifies the norm on the closed disk with the norm on the boundary; therefore the restriction is isometric and injective.

\noindent\textbf{How this problem is used later.} Provides an isometric and unital inclusion in $C(\mathbb T)$ for Problem 1.12(c), where the same function will have different inverses available depending on the algebra.

\subsection*{1.12(c) The same element changes spectrum when changing algebra}

\noindent\textbf{Statement.} For $h(z)=z$, prove $\spec_{A(\mathbb D)}(h)=\overline{\mathbb D}$ and $\spec_{C(\mathbb T)}(h|_{\mathbb T})=\mathbb T$.

\begin{proof}[Solution]
If $|\lambda|>1$, function $(\lambda-z)^{-1}$ is continuous in $\overline{\mathbb D}$ and holomorphic
on $\mathbb D$, so $\lambda-h$ is invertible in $A(\mathbb D)$. If $|\lambda|<1$,
$\lambda-z$ vanishes at $z=\lambda$; if $|\lambda|=1$, any putative inverse would agree on
$\mathbb D$ with $(\lambda-z)^{-1}$, which becomes unbounded when approaching $\lambda$. Therefore
\[
 \spec_{A(\mathbb D)}(h)=\overline{\mathbb D}.
\]

In $C(\mathbb T)$, $\lambda-h$ is invertible exactly when it does not vanish on $\mathbb T$; hence
\[
 \spec_{C(\mathbb T)}(h)=\mathbb T.
\]
\end{proof}

\proofdiagrambranch
{$h(z)=z$ in $A(\mathbb D)$}
{$h|_{\mathbb T}$ in $C(\mathbb T)$}
{Invertibility of $\lambda-h$ depends on whether its reciprocal belongs to the algebra under consideration}
{Invertibility fails in $A(\mathbb D)$ for $\lambda\in\overline{\mathbb D}$; in $C(\mathbb T)$ it fails only on $\mathbb T$}
{$\spec_{A(\mathbb D)}(h)=\overline{\mathbb D}$, $\spec_{C(\mathbb T)}(h)=\mathbb T$}

\begin{cthejemplo}{(The inverse depends on the ambient algebra)}
To compare the two spectra of Problem 1.12(c), let's take $h(z)=z$ in $A(\mathbb D)$ and
its restriction in $C(\mathbb T)$.
In the first algebra, $0-h$ has no holomorphic inverse because $1/z$ has a pole at
$0$; in the second, $h^{-1}(z)=\overline z$ is continuous over $\mathbb T$.
So $0\in\spec_{A(\mathbb D)}(h)$ and $0\notin\spec_{C(\mathbb T)}(h|_{\mathbb T})$.
The same boundary function has an inverse only in the larger algebra. Additionally,
$\overline z$ does not come from $A(\mathbb D)$: if a function $g$ had those boundary values, $zg(z)=1$ in $\mathbb T$ and the maximum modulus principle would give $zg(z)=1$ in all
$\mathbb D$, a contradiction at $z=0$. That is why the inclusion is not of $C^*$-subalgebras.
\end{cthejemplo}

\noindent\textbf{Dependencies.} Spectrum definition and inclusion of Problem 1.12(b) are used; Problem 1.10 calculates the spectrum at $C(\mathbb T)$. In $A(\mathbb D)$ the reciprocal of $\lambda-z$ must also be holomorphic on the disk.

\noindent\textbf{How this problem is used later.} It is the counterexample that motivates the $C^*$-subalgebra hypothesis in spectral permanence from Section 2: a closed and unital Banach subalgebra can change the spectrum.
